\documentclass[a4paper,12pt]{ltjsarticle}
\usepackage[top=25mm,bottom=25mm,left=25mm,right=25mm]{geometry}
\usepackage{booktabs}
\usepackage{array}
\usepackage{longtable}
\usepackage{enumitem}
\usepackage{hyperref}
\usepackage{titlesec}
\usepackage{xcolor}
\usepackage{graphicx}
\usepackage{listings}
\usepackage{luatexja} 
\usepackage{tabularx}

\setlist[itemize]{leftmargin=2em}
\setlist[itemize,2]{leftmargin=2em}
\setlist[enumerate]{leftmargin=2em}
\setlist[enumerate,2]{leftmargin=2em}

\newcolumntype{L}[1]{>{\raggedright\arraybackslash}p{#1}}
\newcolumntype{C}[1]{>{\centering\arraybackslash}p{#1}}

% セクション見出しのスタイル
\titleformat{\section}{\large\bfseries}{第\thesection 節}{1em}{}[\titlerule]
\titleformat{\subsection}{\normalsize\bfseries}{\thesubsection}{1em}{}

% ---- listings 共通設定 ----
\lstset{
  basicstyle      = \ttfamily\small,
  keywordstyle    = \color[HTML]{0000CC}\bfseries,
  commentstyle    = \color[HTML]{008000}\itshape,
  stringstyle     = \color[HTML]{CC0000},
  numberstyle     = \ttfamily\footnotesize\color{gray},
  numbers         = left,          % 行番号を左に
  numbersep       = 10pt,
  stepnumber      = 1,
  frame           = single,        % 枠線
  framesep        = 6pt,
  rulecolor       = \color{gray!40},
  backgroundcolor = \color[HTML]{F7F7F7},
  breaklines      = true,          % 長い行を折り返す
  breakatwhitespace = false,
  showstringspaces= false,
  tabsize         = 4,
  captionpos      = b,             % キャプションを下に
  xleftmargin     = 18pt,          % 行番号のはみ出し対策
  % --- 日本語コメント対策(lualatex) ---
  extendedchars   = true,
}
\renewcommand{\lstlistingname}{ソースコード}

\begin{document}

% ヘッダ情報
\begin{center}
  {\LARGE\bfseries 作業ログ}
\end{center}

\vspace{4mm}

\begin{tabular}{ll}
  \textbf{報告者} & 酒井 睦基 \\
  \textbf{報告日} & \today \\
  \textbf{プロジェクト名} & Arduino オーケストラ：きらきら星 輪唱演奏システム \\
\end{tabular}

\vspace{6mm}

%-------------------------------------------------------------------
\section{進捗サマリー（要約）}

\begin{itemize}
  \item \textbf{全体進捗率}：95 [\%]
  \item \textbf{マイルストーン達成状況}：未達成
  \item \textbf{主要成果}：
    \begin{itemize}
      \item サーボモーターと飾り部分の結合が完了
      \item ポテンショメータの変化に合わせて各小節毎にテンポが速/遅くなることを確認した
      \item PLL補正によって同期誤差が改善されていることが確認
    \end{itemize}
  \item \textbf{未達成要項}：
    \begin{itemize}
      \item 結合テスト（ソフトウェア部分とハードウェア部分を結合したテスト）
    \end{itemize}
  \item \textbf{重大課題（影響度：低）}：7/1（水）のMTGでテストを実施する
\end{itemize}

%-------------------------------------------------------------------
\section{作業ログ（詳細トラッキング）}

\begin{longtable}{C{2cm}C{1.6cm}C{1.6cm}C{2.2cm}C{1.2cm}L{2.4cm}L{3cm}}
  \toprule
  \textbf{作業項目} & \textbf{開始日} & \textbf{予定完了日} & \textbf{状態} & \textbf{進捗率} & \textbf{依存関係・ブロッカー} & \textbf{コメント} \\
  \midrule
  \endfirsthead
  \toprule
  \textbf{作業項目} & \textbf{開始日} & \textbf{予定完了日} & \textbf{状態} & \textbf{進捗率} & \textbf{依存関係・ブロッカー} & \textbf{コメント} \\
  \midrule
  \endhead
  \midrule
  \endfoot
  \bottomrule
  \endlastfoot
  ベビーメリー風回転装置制作 & 6/3 & 6/23 & 完了 & 100\% & 飾り部分含めハードウェアが完成した & 計画書の想定では第7〜8週目には終えている予定であったが，材料調達遅延や不十分な設計，集まれる機会が少なかったことから最終週手前となってしまった \\
  ロードセルの反応検証 & 6/10 & 6/23 & 完了 & 100\% & ロードセルの荷重テスト・キャリブレーションによる閾値の変更まで完了した & 小物の重量によってはシステムが稼働しないためロードセルに乗せる対象の検討を行う \\
  センサ入力・BPM変換処理 & 6/3 & 6/23 & 完了 & 100\% & ポテンショメータによる可変テンポを確認した & 可変テンポに合わせてベビーメリー装置の回転速度が変化することを7/1の結合テストで確認する \\
  PLL補正アルゴリズムの実装 & 6/10 & 6/30 & 完了 & 100\% & PLL補正によって同期誤差が大幅に改善されたことを確認した & 
\end{longtable}

%-------------------------------------------------------------------
\section{技術成熟度レベルTRLの推移}
%-------------------------------------------------------------------
FBSの各要素ごとに現在どのレベルにいるかを以下に示す．

\begin{longtable}{L{2.5cm}L{4cm}C{2cm}C{2cm}L{3.5cm}}
  \toprule
  \textbf{機能} & \textbf{説明} & \textbf{前回TRL} & \textbf{今回TRL} & \textbf{備考} \\
  \midrule
  \endfirsthead
  \toprule
  \textbf{機能} & \textbf{説明} & \textbf{前回TRL} & \textbf{今回TRL} & \textbf{備考} \\
  \midrule
  \endhead
  \midrule
  \endfoot
  \bottomrule
  \endlastfoot
  I2C同期（F1） & マスタ→スレーブ同期通信 & 4 & 4 & マスタ1台，スレーブ4台の接続を確認できた \\
  音色合成（F4） & Processing加算合成 & 4 & 4 & 4種類の楽器音を判別できるレベルで確認できた \\
  PLL補正（F1.3） & 位相同期制御 & 1 & 4 & PLL補正の有無によって同期誤差が大きく変化することが確認できた \\
  エンタメ機能（F5） & ロードセル・サーボ制御 & 2 & 4 & ロードセルの荷重検知によってサーボモーターが稼働し，ベビーメリー飾りが回転することが確認できた \\
  可変テンポ（F2） & BPM制御 & 3 & 4 & ポテンショメータによってテンポが各小節毎に変化することが確認できた \\
  共通ライブラリ（F3.1） & PROGMEM＋I2Cフレーム & 4 & 4 & 共通ライブラリの確認と2年生間の共有を確認できた \\
\end{longtable}

\noindent
{\small
\begin{tabular}{ll}
  \textbf{TRL} & \textbf{定義} \\
  1 & 概念レベル：アイディアや原理が確立されており，説明できるレベル \\
  2 & 基本動作レベル：単体で動作確認が可能なレベル \\
  3 & 結合動作レベル：機能を結合し，実際の使用環境に近い環境で動作確認が可能なレベル \\
  4 & 実運用レベル：実際の使用環境で安定的な稼働が確認できるレベル \\
\end{tabular}
}

%-------------------------------------------------------------------
\section{技術性能指標TPMの推移} 
%-------------------------------------------------------------------
MOPの各要素毎に現在の値の程度を以下に示す．

\begin{longtable}{L{2cm}L{3cm}C{2cm}C{2cm}C{2cm}L{2.5cm}}
  \toprule
  \textbf{関連MOP} & \textbf{指標} & \textbf{目標値} & \textbf{前回値} & \textbf{今回値} & \textbf{備考} \\
  \midrule
  \endfirsthead
  \toprule
  \textbf{関連MOP} & \textbf{指標} & \textbf{目標値} & \textbf{前回値} & \textbf{今回値} & \textbf{備考} \\
  \midrule
  \endhead
  \midrule
  \endfoot
  \bottomrule
  \endlastfoot
  MOP1-1 & 楽器間同期誤差（クラリネット） & $\leq 30$\,ms & 1.35 & -0.38 & 実施済 \\
  MOP1-2 & 楽器間同期誤差（フルート） & $\leq 30$\,ms & 3.07 & -0.01 & 実施済 \\
  MOP1-3 & 楽器間同期誤差（オルガン） & $\leq 30$\,ms & -1.21 & -0.19 & 実施済 \\
  MOP1-4 & 楽器間同期誤差（ドラム） & $\leq 30$\,ms & 4.11 & 0.23 & 実施済 \\
  MOP2-1 & I2C SYNCパケットロス率（クラリネット） & 0回 & 0/50 & 0/50 & 実施済 \\
  MOP2-2 & I2C SYNCパケットロス率（フルート） & 0回 & 0/50 & 0/50 & 実施済 \\
  MOP2-3 & I2C SYNCパケットロス率（オルガン） & 0回 & 0/50 & 0/50 & 実施済 \\
  MOP2-4 & I2C SYNCパケットロス率（ドラム） & 0回 & 0/50 & 0/50 & 実施済 \\
  MOP3 & 楽曲識別（中間確認） & チーム内識別 & 識別可能 & 識別可能 & 実施済 \\
  MOP4 & 楽器識別（中間確認） & チーム内識別 & 識別可能 & 識別可能 & 実施済 \\
  MOP5 & リラクゼーション感 & 心地よく聞こえる & --- & --- & 未実施 \\
  MOP6 & BPM変更反映速度 & 1小節以内 & --- & --- & 未計測 \\
  MOP7 & PLL収束速度 & 変化量$\div$10小節以内 & --- & --- & 未計測 \\
\end{longtable}

\subsection{現在のガントチャート}
現在のガントチャートの状態を図\ref{fig:gc}に示す．
灰色の箇所は実装が完了したことを示す．

\begin{figure}[htbp]
  \centering
  \includegraphics[width=\linewidth]{fig/gc_current.png}
  \caption{現在のガントチャート}
  \label{fig:gc}
\end{figure}

\subsection{日毎の作業時間と作業内容}
定期的にGitHubを確認しPRがきていないかどうかを確認する．
PRがきていた場合はファイルの内容をレビューし，問題がない場合はmergeする．
基本的な作業時間はGitHubの確認とコードレビューを含め1時間程度である．
また，ハードウェア設計としてサーボモーターとベビーメリー飾りの接続が完成した．
サーボモーターの回転に合わせて飾りも回転し，軸がブレて倒れることがないことも確認した．
設計とテストに要した時間はおよそ3時間程度である．

%-------------------------------------------------------------------
\section{課題管理}
\begin{itemize}
  \item \textbf{課題}：結合テスト（ソフトウェア部分とハードウェア部分の結合）
    \begin{itemize}
      \item \textbf{発生原因}：ハードウェア部分の作成が今週完成したため
      \item \textbf{影響度}：高
      \item \textbf{優先度}：高
      \item \textbf{対策}：MTGで結合テストを実施する
      \item \textbf{期限}：7/1
      \item \textbf{状態}：対応中
    \end{itemize}
\end{itemize}

% \noindent
% {\small ※影響度＝プロジェクトへのインパクト \quad ※優先度＝対応の緊急性}

%-------------------------------------------------------------------
\section{AI利用状況と検証（再現性重視）}
今週はハードウェア設計とコードチェックのみを行っていたためAIを利用しなかった．
そのため今回の作業ログではこの節の記載を省く．

\subsection{利用内容}
\begin{itemize}
  \item \textbf{利用目的}：
  \item \textbf{入力（プロンプト要旨）}：
  \item \textbf{出力概要}：
\end{itemize}

\subsection{検証方法}

\begin{itemize}
  \item \textbf{検証手法}：
    \begin{itemize}
      \item 
    \end{itemize}
  \item \textbf{参照資料}：
    \begin{itemize}
      \item 
    \end{itemize}
  \item \textbf{検証手順}：
    \begin{enumerate}
      \item 手順1：
    \end{enumerate}
\end{itemize}

\subsection{ファクトチェック結果}
\begin{itemize}
  \item \textbf{一致点}：
    \begin{itemize}
      \item 
    \end{itemize}
  \item \textbf{不一致・疑義}：
    \begin{itemize}
      \item 
    \end{itemize}
  \item \textbf{最終判断}：
  \item \textbf{根拠}：
    \begin{itemize}
      \item 
    \end{itemize}
\end{itemize}

%-------------------------------------------------------------------
\section{次回計画}

\begin{itemize}
  \item \textbf{全員}：
    \begin{itemize}
      \item 結合テスト
      \item リラクゼーション感のテスト
      \item BPM変更反映速度のテスト
      \item PLL収束速度のテスト
    \end{itemize}
  \item \textbf{期限}：7/1（第10週）
\end{itemize}

%-------------------------------------------------------------------
\section{その他}

\begin{itemize}
  \item 結合テストが無事完了した場合は評価指標の測定に移る．
\end{itemize}

%-------------------------------------------------------------------
\section{付録}

\begin{itemize}
  \item \textbf{添付資料}：最終的な開始トリガ部分のコード
\end{itemize}

\begin{lstlisting}[language=C++, caption=開始トリガシステムのコード, label={code:system}]
#include "HX711.h" // 作者がBogdan NeculaのHX711系のライブラリをインストールする必要がある
#include <Servo.h>

// ロードセルの設定
const int DT_PIN  = 3;
const int SCK_PIN = 2;
HX711 scale;

// サーボモーターの設定
Servo servo;
const int SERVO_PIN = 9;
const float THRESHOLD = 50.0;   // この重さ(g)を超えたら回す．実測値に合わせて変更する
int MOVE_US = 1480;   // 回転速度（BPMの変更に合わせて可変されるようにする）
const int STOP_US   = 1500;   // 停止

void setup() {
  Serial.begin(9600);

  scale.begin(DT_PIN, SCK_PIN);
  scale.set_scale(2280);  // キャリブレーションで設定した値，必要があれば変更を加える
  scale.tare();

  // 初期状態は停止させる
  servo.attach(SERVO_PIN);
  servo.writeMicroseconds(STOP_US);

  Serial.println("準備完了");
}

void loop() {
  float weight = scale.get_units(3);  // 3回平均で重さを取得

  // 荷重検知のテスト用のログ
  Serial.print("重さ: ");
  Serial.print(weight, 1);
  Serial.println(" g");

  // ロードセルに小物が置かれているときはサーボモーターを回転させ，小物が取り除かれたら回転を止める
  if (weight > THRESHOLD) {
    servo.writeMicroseconds(MOVE_US);
  } else {
    servo.writeMicroseconds(STOP_US);
  }
}
\end{lstlisting}

\clearpage
\begin{itemize}
  \item \textbf{添付範囲}：すべて
\end{itemize}

\end{document}